\documentclass[11pt,a4paper]{article}

\usepackage[utf8]{inputenc}
\usepackage[T1]{fontenc}
\usepackage{lmodern}
\usepackage[provide=*,spanish]{babel}
\usepackage[a4paper,margin=2.35cm,headheight=15pt]{geometry}
\usepackage{microtype}
\usepackage{booktabs}
\usepackage{tabularx}
\usepackage{longtable}
\usepackage{array}
\usepackage{xcolor}
\usepackage{hyperref}
\usepackage{xurl}
\usepackage{enumitem}
\usepackage{listings}
\usepackage{fancyhdr}
\usepackage{parskip}
\usepackage{needspace}
\usepackage{graphicx}
\usepackage{pdfpages}
\usepackage{pgf-umlsd}

\definecolor{ucaldasink}{HTML}{2D3142}
\definecolor{ucaldasslate}{HTML}{4F5D75}
\definecolor{ucaldasaccent}{HTML}{BF4520}
\definecolor{ucaldasrule}{HTML}{D8DADF}

\hypersetup{
  colorlinks=true,
  linkcolor=ucaldasink,
  urlcolor=ucaldasaccent,
  citecolor=ucaldasink,
  pdftitle={Gestión de Monitores Académicos},
  pdfauthor={Equipo del proyecto}
}

\setlength{\parindent}{0pt}
\setlength{\parskip}{0.55em}
\setlist[itemize]{leftmargin=1.5em,itemsep=0.2em,topsep=0.25em}
\setlist[enumerate]{leftmargin=1.7em,itemsep=0.25em,topsep=0.25em}
\renewcommand{\arraystretch}{1.2}
\newcolumntype{Y}{>{\raggedright\arraybackslash}X}

\lstdefinestyle{comandos}{
  basicstyle=\ttfamily\small,
  breaklines=true,
  breakatwhitespace=false,
  columns=fullflexible,
  keepspaces=true,
  showstringspaces=false,
  frame=single,
  rulecolor=\color{ucaldasrule},
  backgroundcolor=\color{black!3},
  xleftmargin=0.4em,
  xrightmargin=0.4em,
  aboveskip=0.6em,
  belowskip=0.6em,
  literate=
    {á}{{\'a}}1 {é}{{\'e}}1 {í}{{\'i}}1 {ó}{{\'o}}1 {ú}{{\'u}}1
    {Á}{{\'A}}1 {É}{{\'E}}1 {Í}{{\'I}}1 {Ó}{{\'O}}1 {Ú}{{\'U}}1
    {ñ}{{\~n}}1 {Ñ}{{\~N}}1 {ü}{{\"u}}1 {Ü}{{\"U}}1
}
\lstset{style=comandos}

\pagestyle{fancy}
\fancyhf{}
\fancyhead[L]{\small Gestión de Monitores Académicos}
\fancyhead[R]{\small Universidad de Caldas}
\fancyfoot[L]{\small Ingeniería de Software III · 2026-2}
\fancyfoot[R]{\small Página \thepage}
\renewcommand{\headrulewidth}{0.4pt}
\renewcommand{\footrulewidth}{0.4pt}

\newcommand{\portada}[2]{%
  \hypersetup{pageanchor=false}
  \begin{titlepage}
    \centering
    \vspace*{2.5cm}
    {\large Universidad de Caldas\par}
    \vspace{0.5cm}
    {\large Ingeniería de Software III · 2026-2\par}
    \vspace{2.5cm}
    {\Huge\bfseries Gestión de Monitores Académicos\par}
    \vspace{0.7cm}
    {\LARGE #1\par}
    \vspace{0.7cm}
    {\large #2\par}
    \vfill
    {\large Versión 1.0 · 20 de noviembre de 2026\par}
    \vspace{0.4cm}
    {\normalsize Equipo del proyecto\par}
  \end{titlepage}
  \hypersetup{pageanchor=true}
  \pagenumbering{roman}
}

\newcommand{\codigo}[1]{\texttt{#1}}
\newcommand{\archivo}[1]{\path{#1}}
\newcommand{\componente}[1]{\texttt{#1}}


\begin{document}

\portada{Manual técnico}{Arquitectura, operación y mantenimiento}
\tableofcontents
\clearpage
\pagenumbering{arabic}

\section{Propósito y alcance}

Este manual describe cómo está construido, desplegado y verificado el sistema de Gestión de Monitores Académicos. Está dirigido a quienes desarrollan, integran o mantienen el proyecto. El contrato de API sigue siendo la fuente normativa para rutas, datos, validaciones y errores: \archivo{api-contract.md}.

\textbf{Recorrido rápido}
\begin{enumerate}
  \item Levantar los cinco servicios con \codigo{docker compose up -\kern0pt-build -d}.
  \item Comprobar que todos estén saludables con \codigo{docker compose ps}.
  \item Ejecutar \codigo{./scripts/smoke-test.sh} contra \codigo{http://localhost:8080/api}.
\end{enumerate}

\section{Arquitectura}

La aplicación es una arquitectura distribuida, multicapa y heterogénea. Las capas se comunican mediante HTTP y JSON; las operaciones del CRUD pasan por cada componente en el orden indicado. El navegador solo accede al frontend, y únicamente \archivo{python-service} se conecta a PostgreSQL.

\begin{center}
\textbf{Navegador} $\rightarrow$ \textbf{nginx / frontend} $\rightarrow$ \textbf{Java} $\rightarrow$ \textbf{PHP} $\rightarrow$ \textbf{Python} $\rightarrow$ \textbf{PostgreSQL}
\end{center}

El diagrama de arquitectura, con detalles de tecnologías y puertos, está en \href{run:arquitectura.pdf}{\archivo{arquitectura.pdf}}.

\begin{tabularx}{\textwidth}{p{0.19\textwidth}p{0.25\textwidth}p{0.13\textwidth}Y}
\toprule
\textbf{Componente} & \textbf{Tecnología} & \textbf{Puerto} & \textbf{Responsabilidad} \\
\midrule
\codigo{frontend} & HTML, CSS, JavaScript y nginx & 8080 (host) / 80 (contenedor) & Sirve la interfaz y reenvía \codigo{/api/*} a Java, quitando el prefijo \codigo{/api}. \\
\codigo{java-controller} & Java 21, Spring Boot 3.3 & 8081 & Expone la API, normaliza y valida datos, y propaga las respuestas del servicio PHP. \\
\codigo{php-service} & PHP 8.3 y cURL & 8082 & Comprueba rutas y JSON, y reenvía las solicitudes sin cambiar el contenido. \\
\codigo{python-service} & Python 3.12, FastAPI y psycopg & 8083 & Vuelve a validar, aplica reglas de negocio y ejecuta las operaciones de persistencia. \\
\codigo{database} & PostgreSQL 16 & 5432 & Conserva materias y monitores; el volumen \codigo{db-data} persiste los datos. \\
\bottomrule
\end{tabularx}

\subsection{Responsabilidades y límites}

\begin{itemize}
  \item \textbf{Frontend:} formulario, búsqueda, tabla, confirmación de borrado y presentación de errores. Usa rutas relativas \codigo{/api/...}; no conoce hosts internos ni credenciales.
  \item \textbf{nginx:} sirve los archivos estáticos y actúa como proxy de mismo origen. No implementa reglas de negocio.
  \item \textbf{Java:} normaliza y valida campos antes de enviar la petición a PHP. Los errores de validación se responden sin llamar a la siguiente capa.
  \item \textbf{PHP:} valida método, ruta y sintaxis JSON; entrega a Python método, ruta y cuerpo sin modificarlos.
  \item \textbf{Python:} valida otra vez aunque Java ya lo haya hecho, verifica existencia y duplicados, maneja transacciones y traduce errores de PostgreSQL al contrato.
  \item \textbf{PostgreSQL:} refuerza la integridad con claves primarias, foráneas, restricciones \codigo{UNIQUE} y \codigo{CHECK}.
\end{itemize}

\subsection{Comunicación y ciclo de vida}

Docker Compose crea la red de servicios y controla el orden de inicio mediante \codigo{depends\_on} y comprobaciones de salud: base de datos, Python, PHP, Java y frontend. El estado saludable confirma que el proceso o endpoint responde; no sustituye la prueba completa del flujo CRUD.

Las respuestas del backend regresan por la cadena en sentido inverso. Java y PHP conservan el código HTTP y el cuerpo de la capa siguiente, salvo los errores propios definidos por el contrato. El borrado exitoso responde \codigo{204 No Content}, sin cuerpo.

\section{API y modelo de datos}

\subsection{Operaciones}

\begin{tabularx}{\textwidth}{p{0.14\textwidth}p{0.37\textwidth}p{0.16\textwidth}Y}
\toprule
\textbf{Método} & \textbf{Ruta en la API pública} & \textbf{Éxito} & \textbf{Uso} \\
\midrule
\codigo{GET} & \codigo{/api/health} & 200 & Estado de Java a través del frontend. Cada servicio también expone su propio \codigo{/health}. \\
\codigo{GET} & \codigo{/api/materias} & 200 & Lista materias de solo lectura. \\
\codigo{GET} & \codigo{/api/monitores} & 200 & Lista monitores; sin resultados devuelve \codigo{[]}. \\
\codigo{GET} & \codigo{/api/monitores/\{codigo\}} & 200 & Consulta un monitor por código. \\
\codigo{POST} & \codigo{/api/monitores} & 201 & Crea un monitor. \\
\codigo{PUT} & \codigo{/api/monitores/\{codigo\}} & 200 & Reemplaza los campos editables; no cambia el código. \\
\codigo{DELETE} & \codigo{/api/monitores/\{codigo\}} & 204 & Elimina el monitor sin cuerpo de respuesta. \\
\bottomrule
\end{tabularx}

Las rutas internas vistas por Java, PHP y Python no incluyen el prefijo \codigo{/api}. Las solicitudes de creación y actualización llevan JSON con \codigo{codigo}, \codigo{nombre}, \codigo{correo} y \codigo{materia\_id}; en \codigo{PUT}, el código se toma de la ruta y los otros tres campos son obligatorios. Las respuestas de monitor agregan \codigo{materia\_nombre}.

\subsection{Persistencia}

El esquema se inicializa en \archivo{database/init.sql}. \codigo{materias.id} es la clave referenciada desde \codigo{monitores.materia\_id}; borrar una materia en uso está restringido. El código del monitor es su clave primaria. El correo no es único. Se crea además un índice por \codigo{materia\_id}.

La inicialización automática solo se ejecuta cuando el directorio de datos está vacío. El script carga ocho materias y siete monitores de demostración. El volumen \codigo{db-data} conserva los cambios entre reinicios normales del contenedor.

\section{Validación y manejo de errores}

Java y Python aplican la misma normalización: código sin espacios exteriores y en mayúsculas; nombre sin espacios exteriores y con espacios internos compactados; correo sin espacios exteriores y en minúsculas. La validación responde solo el primer error en un orden determinista.

\begin{tabularx}{\textwidth}{p{0.2\textwidth}Y}
\toprule
\textbf{Dato} & \textbf{Regla principal} \\
\midrule
Código & Obligatorio al crear; de 1 a 20 caracteres: letras ASCII, números o guion. \\
Nombre & Obligatorio; entre 3 y 150 caracteres. \\
Correo & Obligatorio; formato con usuario, arroba y dominio con punto; máximo 150 caracteres. \\
Materia & Entero positivo que corresponda a una fila existente de \codigo{materias}. \\
\bottomrule
\end{tabularx}

Los errores tienen formato \codigo{\{"error": "mensaje"\}}. Estados principales: \codigo{400} para datos inválidos o materia inexistente, \codigo{404} para monitor o ruta inexistente, \codigo{405} para método no permitido, \codigo{409} para código duplicado, \codigo{502} ante indisponibilidad de un servicio descendente, \codigo{503} ante indisponibilidad de PostgreSQL y \codigo{500} para errores no previstos. El detalle técnico se registra en logs, no se devuelve al usuario.

Para el catálogo completo de mensajes exactos y el orden de validación, consultar las secciones 6 a 8 de \archivo{api-contract.md}.

\section{Configuración y despliegue}

\subsection{Requisitos}

Se necesita Docker Desktop con Docker Compose v2 en ejecución y los puertos locales disponibles. No es necesario instalar Java, PHP, Python ni PostgreSQL en el equipo para ejecutar el sistema integrado.

\subsection{Inicio y detención}

\begin{lstlisting}
git clone https://github.com/CristianGarcia21/monitores-ucaldas.git
cd monitores-ucaldas
docker compose up --build -d
docker compose ps
\end{lstlisting}

La interfaz queda disponible en \url{http://localhost:8080}. Para detener los servicios sin borrar la base:

\begin{lstlisting}
docker compose down
\end{lstlisting}

\textbf{Destructivo:} \codigo{docker compose down -v} elimina también el volumen de PostgreSQL y todos los datos. \codigo{./scripts/reset-db.sh} también reinicializa la base de datos.

\subsection{Variables de entorno}

Las variables permiten ajustar los puertos publicados y los parámetros de conexión. Para configurar el entorno local, copiar \codigo{.env.example} como \codigo{.env} y cambiar los valores necesarios.

\begin{tabularx}{\textwidth}{p{0.3\textwidth}p{0.2\textwidth}Y}
\toprule
\textbf{Variable} & \textbf{Por defecto} & \textbf{Uso} \\
\midrule
\archivo{FRONTEND_PORT} & 8080 & Puerto local del frontend. \\
\archivo{JAVA_PORT} & 8081 & Puerto local del controlador. \\
\archivo{PHP_PORT} & 8082 & Puerto local del intermediario. \\
\archivo{PYTHON_PORT} & 8083 & Puerto local del servicio de persistencia. \\
\archivo{DB_PORT} & 5432 & Puerto local publicado de PostgreSQL. \\
\archivo{POSTGRES_DB} & \codigo{monitores} & Nombre de la base de datos. \\
\archivo{POSTGRES_USER} & \codigo{monitores} & Usuario de PostgreSQL. \\
\archivo{POSTGRES_PASSWORD} & \codigo{monitores} & Contraseña de PostgreSQL. \\
\bottomrule
\end{tabularx}

Los valores de conexión entre servicios se configuran en Compose como \archivo{PHP_SERVICE_URL}, \archivo{PYTHON_SERVICE_URL} y \archivo{DB_HOST}. Los nombres internos de host en Docker no se sustituyen por \codigo{localhost}.

\section{Verificación y pruebas}

\subsection{Prueba de integración}

Con los cinco servicios levantados, ejecutar el smoke test completo:

\begin{lstlisting}
./scripts/smoke-test.sh
\end{lstlisting}

En Windows PowerShell:

\begin{lstlisting}
.\scripts\smoke-test.ps1
\end{lstlisting}

El script recorre operaciones CRUD contra \codigo{http://localhost:8080/api}, valida estados y mensajes del contrato, y elimina los registros de prueba al finalizar.

\subsection{Pruebas por servicio}

\begin{itemize}
  \item \textbf{Python:} desde \archivo{python-service}, instalar \archivo{requirements-dev.txt} y ejecutar \codigo{pytest}. Las pruebas que requieren datos necesitan la base activa; las de validación pueden ejecutarse sin ella.
  \item \textbf{Java:} desde la raíz del repositorio, ejecutar:
\begin{lstlisting}
docker run --rm -v "${PWD}/java-controller:/app" \
  -w /app maven:3.9-eclipse-temurin-21 mvn -q test
\end{lstlisting}
  En Git Bash puede requerirse \archivo{MSYS_NO_PATHCONV=1}.
  \item \textbf{PHP:} el README de \archivo{php-service} documenta pruebas con un mock local de Python y peticiones \codigo{curl}.
  \item \textbf{Frontend:} \archivo{frontend/api.js} dispone del modo mock para desarrollo aislado; en la configuración integrada \codigo{USE\_MOCK} es \codigo{false}.
\end{itemize}

\section{Operación y diagnóstico}

\subsection{Salud y logs}

\begin{lstlisting}
docker compose ps
docker compose logs -f java-controller
./scripts/logs.sh python-service
\end{lstlisting}

Cada servicio registra una línea por petición con el patrón \codigo{[servicio] MÉTODO RUTA -> STATUS (ms)}. Los endpoints \codigo{/health} verifican que el proceso respectivo esté vivo, no que todas las dependencias estén disponibles.

\subsection{Problemas frecuentes}

\begin{tabularx}{\textwidth}{p{0.3\textwidth}Y}
\toprule
\textbf{Síntoma} & \textbf{Acción} \\
\midrule
No conecta con Docker & Iniciar Docker Desktop, esperar a que esté listo y repetir el comando. \\
Puerto ocupado & Cambiar el valor correspondiente en \archivo{.env}; si es PostgreSQL, por ejemplo, usar \codigo{DB\_PORT=5433}. \\
Servicio no saludable & Revisar \codigo{docker compose ps} y consultar sus logs con \codigo{docker compose logs -f <servicio>}. \\
La interfaz no refleja cambios en \archivo{database/init.sql} & El script solo corre sobre un volumen vacío. Reinicializar la base elimina los datos existentes. \\
PowerShell bloquea scripts & Ejecutar con \codigo{-ExecutionPolicy Bypass}, según el comando documentado en \archivo{README.md}. \\
\bottomrule
\end{tabularx}

\section{Seguridad y límites conocidos}

La aplicación no implementa autenticación ni autorización. La configuración predeterminada de PostgreSQL es para desarrollo y no debe tratarse como credencial de producción. Mantener el sistema en un entorno controlado y cambiar las credenciales si se expone fuera del equipo local.

El backend evita devolver trazas y ejecuta consultas parametrizadas. El frontend usa un origen único mediante nginx, por lo que no configura CORS. Este proyecto no define copias de seguridad automáticas; el volumen persistente no equivale a un respaldo.

\section{Mapa del repositorio}

\begin{tabularx}{\textwidth}{p{0.36\textwidth}Y}
\toprule
\textbf{Ruta} & \textbf{Contenido} \\
\midrule
\archivo{frontend/} & Interfaz, cliente de API, estilos y configuración nginx. \\
\archivo{java-controller/} & API Java, validación, cliente PHP, excepciones y pruebas. \\
\archivo{php-service/} & Enrutamiento HTTP, validación de JSON y cliente cURL para Python. \\
\archivo{python-service/} & Endpoints, esquemas, repositorio SQL, conexión y pruebas. \\
\archivo{database/init.sql} & Esquema y datos iniciales de PostgreSQL. \\
\archivo{docker-compose.yml} & Construcción, puertos, volúmenes, dependencias y healthchecks. \\
\archivo{scripts/} & Smoke tests y utilidades de logs y reinicio de base de datos. \\
\archivo{docs/api-contract.md} & Contrato normativo de API y datos. \\
\bottomrule
\end{tabularx}

\section{Documentos relacionados}

\begin{itemize}
  \item \archivo{manual-usuario.pdf}: guía de uso de la interfaz.
  \item \archivo{manual-usuario.tex}: fuente LaTeX del manual de usuario.
\item \archivo{arquitectura.pdf}: diagrama de arquitectura.
  \item \archivo{api-contract.md}: contrato detallado de integración.
\end{itemize}

\end{document}
